% Ad Atticum 7.1 (ad-atticum-07-01) --- English translation
% Translated by Claude (Anthropic) from the Latin (Perseus).
% Style guide: STYLE.md.

\ciceroLetterOpener{CICERO ATTICO salutem}

\ciceroSection{1} I had in fact given a letter to Lucius Saufeius, and to you alone --- because there was not time enough for me to write, but I was unwilling that a man so intimate with you should come to you without a letter from me. Still, given the pace at which philosophers travel, I think this one will reach you before that one does. If you have already received that letter, you know that I came to Athens on the day before the Ides of October; that on stepping off the ship at the Piraeus I received your letter from our friend Acastus; that I was thrown into confusion at your having arrived at Rome with a fever, but began to take heart because Acastus reported what I wanted to hear of the recovery of your body; that I shuddered, however, at what your letter brought about Caesar's legions, and pressed you to see that the \textit{[Greek: ambition]} of a certain person you know should do us no harm; and on a matter I had written to you about long ago, on which Turranius had told you otherwise at Brundisium (as I learned from a letter I received from Xeno, an excellent man), I set out briefly why I had not put my brother in charge of the province. That is essentially what was in that letter.

\ciceroSection{2} Now hear the rest. By the fates! Bring all that love with which you have embraced me, and all that good sense of yours which by Hercules I judge to be unique in every field, to bear on this care: that you take thought for my whole position. For I seem to see in prospect a struggle so great --- unless the same god who freed us from the Parthian war more handsomely than we dared to pray for has regard for the commonwealth --- so great a struggle, I say, as has never been. Well, then; this evil is one I share with all. I lay nothing on you to ponder for me there. But that special \textit{[Greek: problem]} that is my own --- please, take it up. Do you see how, with you as my counsellor, I have embraced both men? And how I wish I had listened to you from the beginning, most lovingly as you warned me. \textit{[Greek: But you never persuaded the spirit in my breast]}. At last, however, you did persuade me: to embrace the one because he had deserved supremely well of me, the other because he carried so much weight. So I did, and by every act of compliance I brought it about that neither of them held anyone dearer than me.

\ciceroSection{3} For this is what we were reckoning on: that, joined with Pompey, I should never be forced to do wrong against the commonwealth, and that, holding with Caesar, I should never have to fight against Pompey --- so close was their union. Now there hangs over us, as you point out and I see, the bitterest possible contest between them. And each numbers me as his own, unless one of them is dissembling: for Pompey does not doubt it; he is right to judge that what he now thinks about the commonwealth meets with my strong approval. From both of them I have received, at the same time as yours, letters of such a kind that neither seems to value any man alive more than me.

\ciceroSection{4} But what shall I do? I am not asking about the final extremity --- if it should come to a fight in the field, I see that it is better to be defeated with the one than to conquer with the other --- but about what will be done when I arrive: that there shall be no account taken of his absence; that he should disband his army. ``Speak, Marcus Tullius.'' What shall I say? ``Wait, I beg you, until I have met with Atticus''? There is no room for evasion. Against Caesar, then? Where are those tight-clasped right hands? For it was I who helped to make this possible for him, asked by him in person at Ravenna over the tribune Caelius --- by him in person, did I say? It was Gnaeus our friend, too, in that miraculous third consulship of his. If I take any other line, \textit{[Greek: I shall be ashamed]} not before Pompey alone but before \textit{[Greek: the men and women of Troy]}. \textit{[Greek: Polydamas first of all]} ---

\ciceroSection{5} who? You yourself, of course, the one who praises both my deeds and my writings. So I have escaped this blow through the two earlier consulships of the Marcelli, when Caesar's province was on the table; and shall I now fall headlong into the crisis itself? And so, as ``the fool first gives his opinion,'' my strong preference is to manoeuvre something out of the triumph --- to stay outside the city on the most justifiable of pretexts. They will try all the same to elicit a vote from me. You will perhaps laugh at this point. How I wish I were still lingering in my province! It was plainly the thing to do, if all this was hanging over us. Even so, nothing more wretched. For --- to give you this \textit{[Greek: by the way]} --- I want you to know: all those famous first moves, which you yourself were lifting to the skies in your letters, were \textit{[Greek: a veneer]}.

\ciceroSection{6} How not easy a thing virtue is, and how very difficult the long-drawn pretence of it! For when I thought it right and creditable, out of the yearly allowance which had been decreed me, to leave a yearly stipend to Gaius Coelius the quaestor and pay back into the treasury about a million sesterces, our company groaned: thinking the whole of it ought to be distributed to themselves, so that I should turn out a better friend to the treasuries of the Phrygians and the Cilicians than to our own. But they did not move me; for my good name weighed with me most of all, and yet there was no honour that could be done to anyone that I left out. But let this serve, as Thucydides says, as \textit{[Greek: a digression]} --- and not a useless one.

\ciceroSection{7} You, then, will think about my situation: first, by what art we may keep the goodwill of Caesar; next, about the triumph itself --- which I see, unless the times of the commonwealth get in the way, is \textit{[Greek: easily within reach]}. I judge so both from my friends' letters and from the thanksgiving. The man who voted against it voted for more than if he had voted against all the triumphs ever. Then to him assented one friend of mine, Favonius, and one enemy, Hirrus. Cato, on the other hand, both took part in the writing-up and sent me the most agreeable letter about his own vote. And yet, even as Caesar congratulates me on the thanksgiving, he celebrates a triumph over Cato's vote --- not writing what Cato actually said but only that he had not voted me the thanksgiving.

\ciceroSection{8} I come back to Hirrus. You had begun to bring him round to me; finish the work. You have Scrofa, you have Silius; I have written to them already, and to Hirrus himself. For he had said to them that he could conveniently have blocked the matter, but had been unwilling: that, however, he had assented to Cato, my dearest friend, when he had given the most honourable opinion of me; and that I had sent no letter to him, while sending to everyone else. He spoke the truth. To him alone, and to Crassipes, I had not written.

\ciceroSection{9} And so much for the public business; let us come back home. I want to be quit of that man. He is a sheer \textit{[Greek: schemer]} --- a Lartidius to the life. \textit{[Greek: But]} as for the rest, let us sort it out --- this first, in that it has added care to my grief --- but, all the same, this Precius business, whatever it is, I do not want mixed in with those accounts of mine that he handles. I have written to Terentia and to the man himself: that whatever I can in cash, towards the apparatus of the triumph I hope for, I shall be sending back to you. So I think it will be \textit{[Greek: beyond reproach]} --- but as you please. Take up this care also: by what means we are to bring it off. You yourself indicated something in certain letters sent from Epirus, or perhaps from Athens; and in that I will second you.
